\subsection{Sea Level and Sea Surface Height Equations}\label{sec:sea_level_equations}

This derivation is not a new result, but is included here for completeness
and to clarify the assumptions and approximations made in deriving the
equations for sea level change (SLC) and sea surface height change (SSHC)
in response to ice mass redistribution. The derivation follows the framework
laid out in \textcite{al-attarReciprocitySensitivityKernels2023}.

Throughout, we work within a first-order perturbation theory, assuming that
changes in sea level, displacement, and gravitational potential are all small
relative to the background state. We further assume that the shoreline
geometry is fixed, meaning the ocean function, $C$, does not evolve as sea
level changes; this is what renders the fingerprint problem linear and is
appropriate for the small-amplitude changes associated with modern ice melt.
Finally, we assume that the elastic response of the Earth dominates over
viscous relaxation, which is justified for the timescales relevant to
present-day ice mass flux.

Sea level (SL) is defined upon the equipotential surface such that
%
\begin{equation}
    \Phi(\mathbf{x} + \mathrm{SL}\,\hat{\mathbf{r}}) = \Phi_0,
    \label{eq:sl_definition}
\end{equation}
%
where $\Phi$ is the gravitational potential, $\mathbf{x}$ is the position
vector, $\hat{\mathbf{r}}$ is the unit vector in the radial direction, and
$\Phi_0$ is a constant representing the potential at the sea surface. This
means that the sea surface must adapt to changes in the gravitational and
centrifugal potential to maintain an equipotential surface.

When perturbed by an ice thickness change $\Delta T_I$, four things occur
simultaneously:
%
\begin{itemize}
    \item the gravitational potential itself changes, by $\phi$;
    \item the centrifugal potential changes due to perturbations in
          the Earth's rotation vector, by $\psi$;
    \item the solid Earth surface, $\mathbf{x}$, moves radially by
          the displacement vector $\mathbf{u}$; and
    \item the value of the potential at the ocean surface shifts,
          from $\Phi_0$ to $\Phi_0 + \Phi_g$.
\end{itemize}
%
This results in the SL changing, from SL to $\mathrm{SL} + \Delta\mathrm{SL}$,
in the form of
%
\begin{equation}
    (\Phi + \phi + \psi)\!\left(
        \mathbf{x} + \mathbf{u} + (\mathrm{SL} + \Delta\mathrm{SL})\,\hat{\mathbf{r}}
    \right) = \Phi_0 + \Phi_g,
    \label{eq:sl_perturbed}
\end{equation}
%
to which we apply a first-order Taylor expansion of the left-hand side to get
%
\begin{equation}
    \Phi(\mathbf{x} + \mathrm{SL}\,\hat{\mathbf{r}})
    + \phi + \psi
    + \nabla\Phi \cdot \left(\mathbf{u} + \Delta\mathrm{SL}\,\hat{\mathbf{r}}\right)
    = \Phi_0 + \Phi_g.
    \label{eq:sl_taylor}
\end{equation}
%
This equation can be simplified using:
%
\begin{itemize}
    \item using the relation in \autoref{eq:sl_definition} we can cancel the
          terms involving $\Phi$ and $\Phi_0$;
    \item $\nabla\Phi$ represents the gravitational acceleration vector with
          its radial magnitude as $g$, therefore
          $\nabla\Phi \cdot \hat{\mathbf{r}} = g$; and
    \item the term $\mathbf{u} \cdot \nabla\Phi$ represents the potential
          change due to the vertical motion of the solid surface;
\end{itemize}
%
to give
%
\begin{equation}
    \phi + \psi + \mathbf{u} \cdot \nabla\Phi + g\,\Delta\mathrm{SL} = \Phi_g,
    \label{eq:sl_simplified}
\end{equation}
%
which we can rearrange to get the change in sea level
%
\begin{equation}
    \Delta\mathrm{SL} = -\frac{1}{g}
        \left(\mathbf{u} \cdot \nabla\Phi + \phi + \psi\right)
        + \frac{\Phi_g}{g}.
    \label{eq:slc}
\end{equation}

When calculating SSH, we need to account for the fact that SSH is the
height of the sea surface above a reference ellipsoid, meaning we must
add back the radial displacement of the solid surface. We additionally
choose to remove the centrifugal contribution $\psi$ from the SSH change,
following \textcite{al-attarReciprocitySensitivityKernels2023}; this is a convention motivated by the fact
that satellite altimetry measures height above a reference ellipsoid fixed
to the mean rotational state, so that isolating $\phi$ gives a cleaner
measure of the mass redistribution signal. Whether and how this should be
done is a subtle point discussed in detail in \textcite{tamisieaOngoingGlacialIsostatic2011} and
\textcite{lickleyBiasEstimatesGlobal2018}. Under this convention, SSHC is related to
SLC by
%
\begin{equation}
    \Delta\mathrm{SSH} = \Delta\mathrm{SL}
        + \hat{\mathbf{r}} \cdot \mathbf{u}
        + \frac{\psi}{g}.
    \label{eq:ssh_from_sl}
\end{equation}
%
Since $\nabla\Phi$ points approximately radially with magnitude $g$ at the
Earth's surface, we make the approximation
$\mathbf{u} \cdot \nabla\Phi \approx g\,(\hat{\mathbf{r}} \cdot \mathbf{u})$,
which neglects the contribution of lateral surface displacement to the
potential change. Substituting \autoref{eq:slc} into \autoref{eq:ssh_from_sl}:
%
\begin{equation}
    \Delta\mathrm{SSH} =
    \underbrace{
        -\frac{1}{g}\!\left(
            g\,(\hat{\mathbf{r}} \cdot \mathbf{u}) + \phi + \psi
        \right)
        + \frac{\Phi_g}{g}
    }_{\Delta\mathrm{SL}}
    + \hat{\mathbf{r}} \cdot \mathbf{u}
    + \frac{\psi}{g},
    \label{eq:ssh_expanded}
\end{equation}
%
in which the $\hat{\mathbf{r}} \cdot \mathbf{u}$ terms cancel and the
$\psi/g$ terms cancel, which simplifies to
%
\begin{equation}
    \Delta\mathrm{SSH} = \frac{\Phi_g}{g} - \frac{\phi}{g}.
    \label{eq:sshc}
\end{equation}
%
This result is physically intuitive: SSH change has spatial structure
through $\phi$, the perturbation to the gravitational potential due to
ice mass redistribution, but unlike SL it has no dependence on solid
surface displacement $\hat{\mathbf{r}} \cdot \mathbf{u}$ or centrifugal
potential perturbation $\psi$.

\subsection{Derivation of the relationship between GMSL approximation and SSHC}\label{sec:point_gmsl_derivation}

The motivation of this section is to derive the mathematical relationship between the GMSL approximation and the SSHC values across the altimetry domain, which is used in the traditional method of estimating GMSL change from ocean altimetry data. By defining weights for each datapoint based on their latitude, we can express the GMSL approximation as a weighted sum of the SSHC values using an averaging operator built into \citetitle{al-attarPygeoinfPythonLibrary2025}.

To find the GMSL approximation with regards to SSHC, from \autoref{eq:gmsl_ssh_integral} and \autoref{eq:gmsl_ssh_alt_integral}, we start by considering the integral of the SSHC across a section of the Earth's surface, which we denote as $\Omega$ which has an area of $|\Omega|$:


\begin{equation}
    \gmslc = \dfrac{1}{|\Omega|}\int_\Omega \sshc(\phi, \lambda) \, dA \label{eq:gmsl_sshc_integral}
\end{equation}

Where $\sshc(\phi, \lambda)$ is the sea surface height change at a given longitude $\phi$ and latidue $\lambda$. With a spherical radius of $R$, the area, $dA$, is non-uniform in size and shrinks towards the poles, and can be expressed as:

\begin{align}
    dA &= (R\cos \lambda d\phi)(R d\lambda) \\
    &= R^2 \cos \lambda \, d\phi \, d\lambda
\end{align}

When we take discrete samples across a uniform grid, with spacings of $\Delta \phi$ and $\Delta \lambda$, we convert \autoref{eq:gmsl_sshc_integral} into a double summation over N points:

\begin{equation}
  \gmslc \approx \dfrac{\sum^N_{i=1}\sshc(\mathbf{x}_i)\cdot(R^2\cos \lambda_i \Delta \phi \Delta \lambda)}{\sum^N_{j=1}(R^2 \cos \lambda_j\Delta \phi\Delta \lambda)}
\end{equation}

Which can simplify to:

\begin{equation}
  \gmslc \approx \sum^N_{i=1}\sshc(\mathbf{x}_i)\cdot\dfrac{\cos \lambda_i}{\sum^N_{j=1}\cos \lambda_j}
\end{equation}

Which for ease within an averaging quadrature scheme, we can define the weights for each width, $w_i$, as a ratio of it's local cosine width to the sum of all the cosine widths of the avaliable data points:

\begin{equation}
    w_i = \dfrac{\cos \lambda_i}{\sum^N_{j=1}\cos \lambda_j}
\end{equation}

Which then leads to the final expression for the GMSL approximation as a weighted sum of the SSHC values across the altimetry domain:

\begin{equation}
    \gmslc \approx \sum^N_{i=1} w_i \sshc(\mathbf{x}_i)
\end{equation}


\subsection{Joint operator factorisation}\label{sec:joint_operator_factorisation}

We have the forward operator matrix:

\begin{equation}
    O = \begin{bmatrix}
        \operatorpointocean \operatorsshc \operatorslc \operatorloadice & \operatorpointocean \operatorsshc \operatorslc \operatorloadfirn & \operatorpointocean(\operatorsshc \operatorslc \operatorloadwater + 1) \\
        \operatorpointtg \operatorslc \operatorloadice & \operatorpointtg \operatorslc \operatorloadfirn & \operatorpointtg(\operatorslc \operatorloadwater + 1) \\
        \operatorpointice & \operatorpointice & 0 \\
        \operatorgrace \operatorslc \operatorloadice & \operatorgrace \operatorslc \operatorloadfirn & \operatorgrace \operatorslc \operatorloadwater
    \end{bmatrix}
\end{equation}

Which maps the model space vector $\mathbf{x} = (\thicknesschangeice,\, \thicknesschangefirn,\, \odtchange)$ to the data space vector $d = (d_{\Delta SSH},\, d_{\Delta TG},\, d_{\Delta ISH},\, d_{\text{GRACE}})$.

This is an inefficient operator, as it requires six independent calls to the fingerprint numerical solver $\operatorslc$, which is the most computationally expensive part of the forward model, due to its numerical implementation.

We can factorise this matrix into three sparse matrices $A$, $B$, and $C$ such that $O = ABC$, where $B$ contains only one call to $\operatorslc$, thus reducing the number of fingerprint calls from six to one. The derivation of this factorisation proceeds in three steps:

\textbf{Step 1: Isolate the fingerprint structure}

Notice that $\operatorslc$ appears applied to $\operatorloadice$, $\operatorloadfirn$, $\operatorloadwater$
in rows 1, 2, and 4, but never in row 3. This suggests introducing an intermediate quantity
$\mathbf{u} = \operatorslc\mathbf{L}\,\mathbf{x}$ where $\mathbf{L}$ collects the load operators,
so that $\operatorslc$ is called only once. Define

\begin{equation}
    C = \begin{bmatrix}
        \operatorloadice & \operatorloadfirn & \operatorloadwater \\
        0 & 0 & 1 \\
        1 & 0 & 0 \\
        0 & 1 & 0
    \end{bmatrix}
\end{equation}

where the bottom three rows simply route the state variables
$(\thicknesschangeice,\, \thicknesschangefirn,\, \odtchange)$ to slots that will be
used by the outer operators. The action of $C$ on $\mathbf{x}$ gives the intermediate vector

\begin{equation}
    C\mathbf{x} = \begin{pmatrix}
        \operatorloadice\,\thicknesschangeice + \operatorloadfirn\,\thicknesschangefirn + \operatorloadwater\,\odtchange \\
        \odtchange \\
        \thicknesschangeice \\
        \thicknesschangefirn
    \end{pmatrix}
\end{equation}

\textbf{Step 2: Apply the fingerprint operator selectively}

We want $\operatorslc$ applied to only the first component of $C\mathbf{x}$, leaving the rest
unchanged. This is exactly what a block-diagonal operator achieves:

\begin{equation}
    B = \begin{bmatrix}
        \operatorslc & 0 & 0 & 0 \\
        0 & 1 & 0 & 0 \\
        0 & 0 & 1 & 0 \\
        0 & 0 & 0 & 1
    \end{bmatrix}
\end{equation}

so that

\begin{equation}
    BC\mathbf{x} = \begin{pmatrix}
        \operatorslc(\operatorloadice\,\thicknesschangeice + \operatorloadfirn\,\thicknesschangefirn + \operatorloadwater\,\odtchange) \\
        \odtchange \\
        \thicknesschangeice \\
        \thicknesschangefirn
    \end{pmatrix}
\end{equation}

Linearity of $\operatorslc$ then distributes the first component into
$\operatorslc\operatorloadice\,\thicknesschangeice + \operatorslc\operatorloadfirn\,\thicknesschangefirn
+ \operatorslc\operatorloadwater\,\odtchange$, which is precisely the fingerprint contribution
needed in rows 1, 2, and 4 of $M$.

\textbf{Step 3: Recover the original rows by projection}

It remains to map the four-component vector $BC\mathbf{x}$ back to the four observation types.
Reading off what each row of $M$ requires:
\begin{itemize}
    \item Row 1 needs $\operatorpointocean \operatorsshc$ acting on the fingerprint result, plus $\operatorpointocean$ acting on $\odtchange$ directly.
    \item Row 2 needs $\operatorpointtg$ acting on the fingerprint result, plus $\operatorpointtg$ acting on $\odtchange$ directly.
    \item Row 3 needs $\operatorpointice$ acting on $\thicknesschangeice$ and $\thicknesschangefirn$, with no fingerprint contribution.
    \item Row 4 needs $\operatorgrace$ acting on the fingerprint result only. Unlike rows 1 and 2, there is no direct additive contribution from $\odtchange$, since $\operatorgrace$ extracts the gravitational potential perturbation $\phi$ from the fingerprint response space, which is a globally defined field insensitive to the geometric height redistribution of ODT.
\end{itemize}

This gives

\begin{equation}
    A = \begin{bmatrix}
        \operatorpointocean \operatorsshc & \operatorpointocean & 0 & 0 \\
        \operatorpointtg & \operatorpointtg & 0 & 0 \\
        0 & 0 & \operatorpointice & \operatorpointice \\
        \operatorgrace & 0 & 0 & 0
    \end{bmatrix}
\end{equation}

\textbf{Conclusion}

Combining the three steps, $O = ABC$, where each factor has a clear role: $C$ computes the total load and routes state variables, $B$ applies $\operatorslc$ exactly once to the combined load, and $A$ projects the result onto the four observation spaces.
